% Part: normal-modal-logic
% Chapter: frame-correspondence
% Section: equivalence-S5

\documentclass[../../../include/open-logic-section]{subfiles}

\begin{document}

\olfileid{nml}{frd}{es5}

\olsection{তুল্যতা সম্পর্ক ও \Log{S5}}

মোডাল যুক্তি \Log{S5} হলো সব সর্বজনীন কাঠামোতে বৈধ !!{formula}s-এর সমষ্টি। সর্বজনীন কাঠামোয় প্রত্যেক জগৎ থেকে নিজের জগৎসহ প্রত্যেক জগৎ অভিগম্য। এই অবস্থায় $\Box$ আবশ্যিকতা এবং $\Diamond$ সম্ভাব্যতার সঙ্গে মেলে: \emph{প্রত্যেক} জগতে $!A$ সত্য হলে $\Box !A$ সত্য, আর \emph{কোনো একটি} জগতে $!A$ সত্য হলে $\Diamond !A$ সত্য। দেখা যায়, সব প্রতিবিম্বী, প্রতিসম ও পরিযায়ী কাঠামোতে, অর্থাৎ সব \emph{তুল্যতা সম্পর্কের} উপর, বৈধ !!{formula}s-এর সমষ্টি হিসেবেও \Log{S5}-কে চিহ্নিত করা যায়।

\begin{defn}
  $W$-এর উপর কোনো দ্বিপদী সম্পর্ক $R$ \emph{তুল্যতা সম্পর্ক} তখনই এবং কেবল তখনই, যখন সেটি প্রতিবিম্বী, প্রতিসম ও পরিযায়ী। $W$-এর উপর সম্পর্ক $R$ \emph{সর্বজনীন} তখনই এবং কেবল তখনই, যখন সব $u,v \in
  W$-এর জন্য $Ruv$।
\end{defn}

\Ax{T}, \Ax{B} ও \Ax{4} যথাক্রমে প্রতিবিম্বী, প্রতিসম ও পরিযায়ী কাঠামো চিহ্নিত করে। তাই যে কাঠামোগুলিতে অভিগম্যতা সম্পর্কটি তুল্যতা সম্পর্ক, সেগুলি ঠিক সেই কাঠামো, যেখানে তিনটি !!{formula}s-ই বৈধ। তুল্যতা সম্পর্ককে অন্য !!{formula}s-এর সমন্বয় দিয়েও চিহ্নিত করা যায়; কারণ তুল্যতা সম্পর্কের সংজ্ঞার শর্তগুলি~$R$-এর অন্য পরিচিত শর্তসমন্বয়ের সমতুল্য।

\begin{prop}\ollabel{prop:equivalences}
  নিচের শর্তগুলি সমতুল্য:
  \begin{enumerate}
  \item $R$ একটি তুল্যতা সম্পর্ক;
  \item $R$ প্রতিবিম্বী ও ইউক্লিডীয়;
  \item $R$ সিরিয়াল, প্রতিসম ও ইউক্লিডীয়;
  \item $R$ সিরিয়াল, প্রতিসম ও পরিযায়ী।
  \end{enumerate}
\end{prop}

\begin{proof}
  অনুশীলন।
\end{proof}

\begin{prob}
  নিম্নলিখিত বিষয়গুলি দেখিয়ে \olref[nml][frd][es5]{prop:equivalences} প্রমাণ করুন:
  \begin{enumerate}
  \item $R$ প্রতিসম ও পরিযায়ী হলে ইউক্লিডীয়।
  \item $R$ প্রতিবিম্বী হলে সিরিয়াল।
  \item $R$ প্রতিবিম্বী ও ইউক্লিডীয় হলে প্রতিসম।
  \item $R$ প্রতিসম ও ইউক্লিডীয় হলে পরিযায়ী।
  \item $R$ সিরিয়াল, প্রতিসম ও পরিযায়ী হলে প্রতিবিম্বী।
  \end{enumerate}
  শর্তগুলি সমতুল্য প্রমাণে এগুলিই কেন যথেষ্ট, তা ব্যাখ্যা করুন।
\end{prob}

\olref{prop:equivalences} হলো \olref[prf][prs]{prop:S5}-এর অর্থতাত্ত্বিক অনুরূপ: এটি $R$ তুল্যতা সম্পর্ক এমন কাঠামোর মোডাল যুক্তিকে (প্রচলিত নাম \Log{S5}) সমতুল্য শর্তে চিহ্নিত করে।

সর্বজনীন ও তুল্যতা সম্পর্কের মধ্যে কী সম্পর্ক? প্রত্যেক সর্বজনীন সম্পর্ক তুল্যতা সম্পর্ক হলেও, প্রত্যেক তুল্যতা সম্পর্ক সর্বজনীন নয়। কিন্তু সব সর্বজনীন সম্পর্কের উপর বৈধ !!{formula}s ঠিক সেই সূত্রগুলি, যেগুলি সব তুল্যতা সম্পর্কের উপর বৈধ।

\begin{prop}
  ধরা যাক, $R$ একটি তুল্যতা সম্পর্ক। প্রত্যেক $w \in W$-এর \emph{তুল্যতা শ্রেণি}~$w$-এর জন্য হলো $[w] = \{w'\in W :
  Rww'\}$। তাহলে:
  \begin{enumerate}
  \item $w \in [w]$;
  \item প্রত্যেক তুল্যতা শ্রেণি $[w]$-এর উপর $R$ সর্বজনীন;
  \item তুল্যতা শ্রেণিগুলি $W$-কে পরস্পর অসংযুক্ত এবং সম্মিলিতভাবে নিঃশেষকারী উপসমষ্টিতে বিভক্ত করে।
  \end{enumerate}
\end{prop}

\begin{prop}\ollabel{prop:S5=univ}
  !!^a{formula}~$!A$ এমন সব কাঠামো $\mModel{F} =\tuple{W,R}$-এ বৈধ, যেখানে $R$ তুল্যতা সম্পর্ক, তখনই এবং কেবল তখনই, যখন এটি এমন সব কাঠামো $\mModel{F} =\tuple{W,R}$-এ বৈধ, যেখানে $R$ সর্বজনীন। অতএব, সর্বজনীন কাঠামোর যুক্তিই \Log{S5}।
\end{prop}

\begin{proof}
  $W$-এর উপর সর্বজনীন সম্পর্ক~$R$ যে তুল্যতা সম্পর্ক, তা সরাসরি যাচাই করা যায়। তাই $R$ তুল্যতা সম্পর্ক এমন সব কাঠামোতে $!A$ বৈধ হলে সব সর্বজনীন কাঠামোতেও তা বৈধ। অন্য দিকের জন্য বিপরীত-অনুবর্তন ব্যবহার করি: ধরা যাক, $!B$ এমন একটি !!a{formula}, যা $W$-এর উপর তুল্যতা সম্পর্ক $R$-বিশিষ্ট কাঠামো $\tuple{W,R}$-এর উপর প্রতিষ্ঠিত মডেল $\mModel{M} = \tuple{W, R, V}$-এর কোনো জগৎ $w$-তে ব্যর্থ। অর্থাৎ $\mSat/{M}{!B}[w]$। এবার মডেল $\mModel{M}' = \tuple{W',
    R', V'}$ এভাবে সংজ্ঞায়িত করি:
  \begin{enumerate}
  \item $W' = [w]$;
  \item $W'$-এর উপর $R'$ সর্বজনীন;
  \item $V'(p) = V(p) \cap W'$।
  \end{enumerate}
  (অর্থাৎ, $\mModel{M}'$-এর জগৎসমষ্টি $W'$-কে \olref{fig:partition}-এ ছায়াযুক্ত অংশে দেখানো হয়েছে।) $W'$-এর উপর $R$ ও $R'$ যে একই, তা সহজেই বোঝা যায়। এরপর !!{formula}s-এর উপর আবেশে দেখানো যায়, সব $w' \in W'$-এর জন্য: $\mSat{M'}{!A}[w']$ তখনই এবং কেবল তখনই, যখন প্রত্যেক $!A$-এর জন্য $\mSat{M}{!A}[w']$ (কারণ $W'
  \subseteq W$)। বিশেষত $\mSat/{M'}{!B}[w]$; ফলে সর্বজনীন কাঠামোর উপর প্রতিষ্ঠিত একটি মডেলেও $!B$ ব্যর্থ।
\end{proof}

\begin{figure}[t]
  \centering
  \begin{tikzpicture}[node distance=2cm, auto, thick]
    \clip [rounded corners] (0,0) -- (6,0) -- (6,4) -- (0,4) --  cycle;
    \begin{scope}
      \clip (0,0) -- (2,0)  .. controls (1.5,1.5) and (4.5,1.5)
      .. (4,4) -- (0,4) -- (0,0) -- cycle;
      \filldraw[fill=gray!40] (0,4) -- (0,2) .. controls (1.5,1.5)
      and (4.5,1.5) .. (4.5,0) -- (6,0) -- (6,4) -- (0,4) --  cycle;
    \end{scope}
    \draw [name path=line1] (2,0) .. controls (1.5,1.5) and (4.5,1.5) .. (4,4);
    \draw [name path=line2] (0,2)  .. controls (1.5,1.5) and (4.5,1.5) .. (4.5,0);
    \path [name intersections={of=line1 and line2}];
    \draw [rounded corners, use as bounding box, very thick] (0,0)
    -- (6,0) -- (6,4) -- (0,4) --  cycle;
    \node at (2,3) {$[w]$} ;
    \node at (1,1) {$[u]$} ;
    \node at (3,0.75) {$[v]$} ;
    \node at (4.75,2) {$[z]$} ;
  \end{tikzpicture}
  \caption{তুল্যতা শ্রেণিতে $W$-এর বিভাজন।}
  \ollabel{fig:partition}
\end{figure}

\end{document}
